\section{Truco Paulista Rules}
\label{app:rules}

This appendix provides the complete rules of Truco Paulista as implemented in TrucoBench.

\subsection{Deck and Cards}

A 40-card deck is used (standard 52-card deck minus all 8s, 9s, and 10s). The four suits are: Ouros ($\diamondsuit$), Espadas ($\spadesuit$), Copas ($\heartsuit$), Paus ($\clubsuit$).

Card strength from weakest to strongest:
$$4 < 5 < 6 < 7 < Q < J < K < A < 2 < 3$$

Suits do not break ties for non-manilha cards.

\subsection{Manilhas}

At the start of each round, one card is flipped as the \emph{vira}. The four cards of the next rank in the cycle become \emph{manilhas}---the strongest cards for that round.

Manilha suit ordering (weakest to strongest):
$$\diamondsuit < \spadesuit < \heartsuit < \clubsuit \text{ (Zap)}$$

The Zap ($\clubsuit$ manilha) is the single strongest card in the game.

\subsection{Round Structure}

Each round is best-of-3 tricks. Each trick: both players play one card, higher card wins. Equal-strength cards result in a draw.

Draw resolution:
\begin{itemize}
    \item Trick 1 draws $\to$ trick 2 winner takes the round
    \item Trick 1 has a winner, trick 2 draws $\to$ trick 1 winner takes the round
    \item Trick 3 draws $\to$ trick 1 winner takes the round
    \item All three tricks draw $\to$ first player of the round wins
\end{itemize}

\subsection{Escalation}

At any time during their turn, a player may call \emph{Truco} instead of playing a card:

\begin{center}
\begin{tabular}{lc}
\toprule
Level & Points \\
\midrule
Normal & 1 \\
Truco & 3 \\
Seis & 6 \\
Nove & 9 \\
Doze & 12 \\
\bottomrule
\end{tabular}
\end{center}

The opponent responds with Accept, Raise, or Fold. On fold, the caller wins the previous level's points.

\subsection{M\~{a}o de Onze}

When a team reaches exactly 11 points:
\begin{itemize}
    \item They see their cards and choose to play (round worth 3 points) or fold (opponent gets 1 point).
    \item They cannot call Truco during this round.
\end{itemize}

When both teams are at 11 (\emph{m\~{a}o de ferro}): the round is worth 3 points with no escalation allowed.

\subsection{Game End}

The first team to reach 12 points wins.

\section{Prompt Examples}
\label{app:prompts}

The following examples show the three prompt variants for the same game state: Player holds 3$\spadesuit$ (manilha), 2$\spadesuit$, A$\heartsuit$. Vira is 2$\heartsuit$. Opponent has played K$\clubsuit$. Score is 0--0.

\subsection{Minimal Prompt}

\begin{verbatim}
Hand: 3 of Espadas (manilha), 2 of Espadas, A of Copas
Vira: 2 copas
Score: You 0 - Opponent 0
Opp played: K♣
Current stake: 1 point
Actions:
1. PLAY_CARD 0 (3 of Espadas (manilha))
2. PLAY_CARD 1 (2 of Espadas)
3. PLAY_CARD 2 (A of Copas)
4. TRUCO
\end{verbatim}

\subsection{Standard Prompt}

\begin{verbatim}
You are playing Truco Paulista.

## Your hand
- 0: 3 of Espadas (Spades) (manilha)
- 1: 2 of Espadas (Spades)
- 2: A of Copas (Hearts)

## Vira
2 of Copas (Hearts)

## Manilhas (strongest to weakest)
1. 3 of Paus (Clubs) (Zap)
2. 3 of Copas (Hearts)
3. 3 of Espadas (Spades)
4. 3 of Ouros (Diamonds) (weakest)

## Score
You: 0 | Opponent: 0

## Current round (trick 1 of 3)
- Trick 1: Opponent plays K♣. Your turn.

## Escalation
Current stake: 1 point

## Legal actions
1. PLAY_CARD 0 (3 of Espadas (manilha))
2. PLAY_CARD 1 (2 of Espadas)
3. PLAY_CARD 2 (A of Copas)
4. TRUCO

Respond in JSON:
{"reasoning": "...", "action": "PLAY_CARD", "card_index": 0}
\end{verbatim}

\subsection{Verbose Prompt}

The verbose variant prepends the standard prompt with a complete rules summary and strategic hints (approximately 200 additional tokens). See the repository for the full template.

\section{Full Results Tables}
\label{app:results}

% TODO: Include detailed results tables after experiments
